\documentclass[12pt]{article}
\usepackage[utf8]{inputenc}
\usepackage[T1]{fontenc}
\usepackage{hyperref}
\title{Takwane Language Resources for K-12 Curriculum Development}
\author{Compiled by Claude Sonnet 5 (claude-sonnet-5) for Andr\'as Kornai}
\date{}
\begin{document}
\maketitle

\section{Introduction}
Takwane (Thakwani; ISO 639-3 tke) is a Bantu language of Zambezia province, Mozambique. Joshua Project reports about 250,000 people; the survey reference gives about 180,000 speakers. Resources found are mainly scripture and a small dictionary.

\section{Primary Educational Resources}
No primers, textbooks or school programmes were found.

\section{Linguistic and Reference Materials}
\begin{itemize}
\item Webonary Takwane dictionary (trilingual, work in progress, extracted from a database based on the Central African Word List of about 2,000 words): \url{https://www.webonary.work/takwane/overview/introduction/?lang=en}.
\end{itemize}

\section{Religious and Community Resources}
\begin{itemize}
\item eBible.org, Takwane New Testament (Wycliffe, 2007--2009, CC BY-ND 4.0): \url{https://ebible.org/bible/details.php?id=tke}. PDFs: \url{https://ebible.org/pdf/tke/tke_all.pdf}.
\item ScriptureEarth, Takwane (New Testament audio and some Old Testament portions): \url{https://www.scriptureearth.org/00eng.php?iso=tke}.
\item Digital Bible Society: \url{https://dbs.org/bibles/TKEWBT}.
\item Joshua Project, Takwane in Mozambique: \url{https://joshuaproject.net/people_groups/19016}.
\end{itemize}

\section{Digital and Audio Resources}
Scripture audio is available through ScriptureEarth (above). No other audio, video or corpus was found.

\section{Curriculum Development Context}
No school use was found. The New Testament and the small dictionary are the only substantial texts located, so a curriculum would first need literacy materials produced with the community.
\end{document}